\section{Field shielding: $B(z)$}
\label{sec:field}

\subsection{The field is a cone, not an array}

In a directionally drilled island field the wellheads are close and the wells
diverge. THUMS wellheads sit on \SI{6}{ft} (\SI{1.83}{\metre}) centres with
\SI{12}{ft} between double rows, and over 1200 wells reach some 6500 acres
with drift angles to \SI{84}{\degree}. Spacing is therefore a function of
depth, and over the storage interval it varies by more than an order of
magnitude.

The governing quantity is the time for neighbouring thermal fields to merge,
$\tau = B^{2}/\alpha$:

\begin{center}
\begin{tabular}{lr}
\toprule
spacing & merge time \\
\midrule
\SI{1.83}{\metre} (wellhead) & \SI{36}{days} \\
\SI{8}{\metre}               & \SI{1.9}{yr}  \\
\SI{70}{\metre} (base of storage) & \SI{145}{yr} \\
\SI{162}{\metre} (reservoir) & \SI{779}{yr}  \\
\bottomrule
\end{tabular}
\end{center}

The shallow section is a single thermal body within a month; the deep section
is a field of isolated wells. Both occur in the same well.

\paragraph{The two effects oppose one another with depth.} Shallow rock is
cold, so the driving difference is largest exactly where the shielding is
strongest; deep rock is hot, so the driving difference is smallest exactly
where the wells stand alone. This partial cancellation is specific to
directionally drilled island fields and does not arise in purpose-drilled
borehole thermal storage, where the spacing is uniform by design. It is, as
far as we are aware, untreated in the literature, and it is a structural
property of repurposing this particular class of infrastructure.

\subsection{Model}

Above the kickoff point the wells are vertical; below it they fan at an
effective half-angle $\theta$, so with
$R_0 = B_0\sqrt{N/\pi}$,
\begin{equation}
R_f(z) = R_0 + \max\!\big(0,\ \text{depth}(z) - z_{\mathrm{kop}}\big)\tan\theta,
\qquad
B(z) = B_0\,\frac{R_f(z)}{R_0}.
\end{equation}
A real directional programme builds angle over a section and then holds it, so
the fan is not a straight cone; the quantity this feeds is a logarithm, and
the difference is well inside the uncertainty in the spacing data itself.

A shielded well cannot lose what an isolated one would. Once the cells have
merged, the only heat leaving crosses the cluster perimeter, and it divides
among $N$ wells. Taking the semi-infinite-slab flux at that boundary,
$q'' = k_r \Delta T/\sqrt{\pi\alpha t}$, over a perimeter $2\pi R_f(z)$ gives
a loss per well of $q'' \, 2\pi R_f(z)/N$, which written as a resistance is
\begin{equation}
\Rper(z) = \frac{N\sqrt{\pi\alpha t}}{2\pi k_r R_f(z)}\,.
\label{eq:Rper}
\end{equation}
This carries no driving temperature and is therefore a true resistance, so it
enters the series chain directly.

\subsection{Combining the two regimes}

The two are combined by taking the larger resistance:
\begin{equation}
R_{\mathrm{form}}(z) = \Rnear + \max\!\big(\Rg,\ \Rper(z)\big).
\label{eq:maxrule}
\end{equation}
This is not a smooth blend and does not pretend to be. A well cannot lose more
heat than if it stood alone, and a field cannot lose more than its perimeter
admits; whichever constraint binds is the one that governs. The crossover is
sharp in reality as well --- it is the moment the cells merge, and $\tau =
B^2/\alpha$ above shows that moment to be well defined.

\begin{table}[h]
\centering
\begin{tabular}{rrrr}
\toprule
depth, \si{\metre} & $B(z)$, \si{\metre} & $R_{\mathrm{form}}$, \si{\metre\kelvin\per\watt}
 & shielding vs isolated \\
\midrule
0    & 1.8  & 63.62 & $116\times$ \\
300  & 1.8  & 63.62 & $116\times$ \\
600  & 18.6 & 6.39  & $11.7\times$ \\
1000 & 41.1 & 2.99  & $5.5\times$ \\
1524 & 70.4 & 1.81  & $3.3\times$ \\
\bottomrule
\end{tabular}
\caption{$N = 1000$, $\theta = \SI{45}{\degree}$, $t_{\mathrm{op}} =
\SI{10}{yr}$. The shielding weakens by a factor of 35 down the well but never
stops binding: even at the base, where the wells are \SI{70}{\metre} apart,
the field resistance is more than three times the isolated one.}
\end{table}

